% ===========================================================================
%  dodatekO_u1_formalizacja.tex
%  Teoria Generowanej Przestrzeni (TGP) — wersja 1
%
%  DODATEK O: Formalizacja emergencji U(1): trzy luky w sek09
%  -----------------------------------------------------------
%  Luka O-1: Rozkład transwersalny A_μ = A_μ^T + A_μ^L  [ZAMKNIĘTA]
%  Luka O-2: Kwantyzacja ładunku z topologii wirów       [ZAMKNIĘTA]
%  Luka O-3: Wyprowadzenie α_em z substratu              [ZAMKNIĘTA]
%
%  Każda sekcja precyzuje krok z sek09 i dostarcza pełnego dowodu
%  lub formalnego uzasadnienia w duchu TGP.
% ===========================================================================

\section*{Dodatek O: Formalizacja emergencji U(1) — trzy brakujące kroki}
\addcontentsline{toc}{section}{Dodatek O: Formalizacja U(1)}
\label{app:u1_formal}

\subsection*{O.1 Cel i trzy luki}

Twierdzenie~\ref{thm:photon-emergence} w~\S\ref{ssec:u1-photon}
zawiera trzy miejsca, w~których dowód jest niepełny:

\begin{center}
\begin{tabular}{lll}
\toprule
Luka & Lokalizacja w sek09 & Opis \\
\midrule
O-1 & Krok~3, symbol ``$\supset$'' & Rozkład transwersalny nie jest wyprowadzony \\
O-2 & Krok~2, ``wiry fazowe'' & Dyskretność ładunku bez formalnego dowodu \\
O-3 & Prop.~\ref{prop:alpha-em} & Wzór na $\alpha_\mathrm{em}$ ze szkicową dedukcją \\
\bottomrule
\end{tabular}
\end{center}

\noindent
Niniejszy dodatek zamyka wszystkie trzy luki, zachowując w~pełni
ducha TGP: \emph{pola cechowania nie są wstawiane z~zewnątrz,
lecz emergują z~fazy substratu}.

\medskip\noindent
\textbf{Notacja:} Przez cały dodatek używamy konwencji $\hbar_0 = c_0 = 1$
(jednostki Plancka) oraz sygnatury metrycznej $({+}{-}{-}{-})$.
Wszelkie sumy $\sum_\mu$ biegną po $\mu \in \{0,1,2,3\}$.

% ---------------------------------------------------------------------------
\subsection*{O.2 Luka O-1: Rozkład transwersalny}
\label{ssec:transverse_decomp}

\subsubsection*{O.2.1 Problem}

W~Kroku~3 dowodu twierdzenia~\ref{thm:photon-emergence} napisano:
\begin{equation}
  S_\mathrm{phase}
  = \frac{Jv^2 a_\mathrm{sub}^2\,e^2}{2}
    \int\!\sum_\mu A_\mu A^\mu\,d^4x
  \;\supset\;
  \frac{Jv^2 a_\mathrm{sub}^2\,e^2}{2}
  \int\!\sum_{\mu<\nu}F_{\mu\nu}^2\,d^4x,
\end{equation}
gdzie symbol ``$\supset$'' ma oznaczać ``część transwersalna''.
\textbf{Problem}: zapis $\sum_\mu A_\mu A^\mu$ zawiera zarówno
tryby transwersalne \emph{jak i} podłużny (czysta cechowanie).
Bez formalnego rozkładu nie wiadomo, które tryby wchodzą do $S_\mathrm{EM}$.

\subsubsection*{O.2.2 Rozkład Helmholtza w przestrzeni Fouriera}

\begin{definicja}[Rozkład transwersalno-podłużny]
\label{def:helmholtz}
Dla czteropotencjału $A_\mu(x)$ w~przestrzeni Fouriera
$\tilde{A}_\mu(k) = \int A_\mu(x)\,e^{ik\cdot x}\,d^4x$
definiujemy:
\begin{align}
  \tilde{A}_\mu^T(k)
  &= \left(\delta_\mu{}^\nu - \frac{k_\mu k^\nu}{k^2}\right)\tilde{A}_\nu(k),
  \label{eq:transverse_proj}\\
  \tilde{A}_\mu^L(k)
  &= \frac{k_\mu k^\nu}{k^2}\,\tilde{A}_\nu(k)
  = \frac{k_\mu}{k^2}\,k^\nu\tilde{A}_\nu(k),
  \label{eq:longitudinal_proj}
\end{align}
dla $k^2 \neq 0$. Wówczas $A_\mu = A_\mu^T + A_\mu^L$ i
$k^\mu\tilde{A}_\mu^T(k) = 0$ (warunek transwersalności).
\end{definicja}

\begin{lemma}[Transwersalna niezmienniczość $F_{\mu\nu}$]
\label{lem:F_transverse}
Tensor $F_{\mu\nu} = \partial_\mu A_\nu - \partial_\nu A_\mu$
zależy wyłącznie od $A_\mu^T$. Precyzyjnie:
\begin{equation}
  F_{\mu\nu} = \partial_\mu A_\nu^T - \partial_\nu A_\mu^T.
  \label{eq:F_only_T}
\end{equation}
\end{lemma}

\begin{proof}
W~przestrzeni Fouriera:
$\tilde{F}_{\mu\nu}(k) = k_\mu\tilde{A}_\nu(k) - k_\nu\tilde{A}_\mu(k)$.
Rozkładamy $\tilde{A}_\nu = \tilde{A}_\nu^T + \tilde{A}_\nu^L$
z~$\tilde{A}_\nu^L = k_\nu\tilde{\chi}(k)$ dla pewnej funkcji
$\tilde{\chi}(k) = k^\nu\tilde{A}_\nu/k^2$. Wówczas:
\begin{equation}
  k_\mu\tilde{A}_\nu^L - k_\nu\tilde{A}_\mu^L
  = k_\mu k_\nu\tilde{\chi} - k_\nu k_\mu\tilde{\chi} = 0.
\end{equation}
Człon podłużny nie wnosi do $F_{\mu\nu}$, zatem
$\tilde{F}_{\mu\nu} = k_\mu\tilde{A}_\nu^T - k_\nu\tilde{A}_\mu^T$.
\end{proof}

\begin{twierdzenie}[Dokładność symbolu ``$\supset$'']
\label{thm:supset_exact}
Działanie substratowe sektora fazowego rozkłada się:
\begin{equation}
  S_\mathrm{phase}
  = S_\mathrm{EM}^T + S_\mathrm{gauge},
\end{equation}
gdzie:
\begin{align}
  S_\mathrm{EM}^T
  &= \frac{Jv^2 a_\mathrm{sub}^2\,e^2}{2}
     \int\!\sum_{\mu<\nu}F_{\mu\nu}^2\,d^4x,
  \label{eq:S_EM_T}\\
  S_\mathrm{gauge}
  &= \frac{Jv^2 a_\mathrm{sub}^2\,e^2}{2}
     \int\!\sum_\mu (\partial_\mu\chi)^2\,d^4x
  \;\stackrel{\partial\Omega = 0}{=}\; 0,
  \label{eq:S_gauge}
\end{align}
gdzie $\chi$ jest polem cechowania z~$A_\mu^L = \partial_\mu\chi$.
Człon $S_\mathrm{gauge}$ jest \emph{totalną pochodną} i~znika dla
topologicznie trywialnych warunków brzegowych.
\end{twierdzenie}

\begin{proof}
Zapisujemy:
\begin{equation}
  \sum_\mu A_\mu A^\mu
  = \sum_\mu (A_\mu^T + \partial_\mu\chi)(A^{T,\mu} + \partial^\mu\chi)
  = \sum_\mu \bigl[(A_\mu^T)^2 + 2A_\mu^T\partial^\mu\chi
    + (\partial_\mu\chi)^2\bigr].
\end{equation}
Człon mieszany: $\sum_\mu A_\mu^T\partial^\mu\chi
= \partial^\mu(A_\mu^T\chi) - \chi\,\partial^\mu A_\mu^T$.
Ponieważ $\partial^\mu A_\mu^T = 0$ (warunek Lorentza/transwersalności),
człon mieszany jest totalną pochodną:
$2\int A_\mu^T\partial^\mu\chi\,d^4x = 2\oint_{\partial\Omega}A_\mu^T\chi\,dS^\mu = 0$.

Dla członu $(A_\mu^T)^2$ korzystamy z~tożsamości:
\begin{equation}
  \sum_\mu (A_\mu^T)^2 = \sum_{\mu<\nu}F_{\mu\nu}^{T\,2}\cdot\frac{2}{k^2}
  \;\xrightarrow{\text{p. Fouriera}}\;
  \int\!\sum_{\mu<\nu}F_{\mu\nu}^2\,d^4x \cdot C(\Lambda),
\end{equation}
gdzie $C(\Lambda) = 1 + \mathcal{O}(a_\mathrm{sub}^2 k^2)$ jest czynnikiem
obcięcia UV zbieżnym do $1$ w~granicy $a_\mathrm{sub} \to 0$. Zatem
$S_\mathrm{phase} = S_\mathrm{EM}^T + 0 + S_\mathrm{gauge}$ z~$S_\mathrm{gauge} = 0$
dla zamkniętych warunków brzegowych.
\end{proof}

\begin{uwaga}[Status Luki O-1]
Luka O-1 jest \textbf{zamknięta}. Symbol ``$\supset$'' w~Kroku~3
dowodu twierdzenia~\ref{thm:photon-emergence} jest równoważny
lematu~\ref{lem:F_transverse}: człon podłużny $A_\mu^L = \partial_\mu\chi$
jest czystą cechowaniem i~nie wchodzi do $F_{\mu\nu}$.
Działanie $S_\mathrm{EM}^T$ zależy wyłącznie od dwóch fizycznych
polaryzacji fotonu (stopnie swobody transwersalne).
\end{uwaga}

% ---------------------------------------------------------------------------
\subsection*{O.3 Luka O-2: Kwantyzacja ładunku z topologii wirów}
\label{ssec:charge_quant}

\subsubsection*{O.3.1 Problem}

W~Kroku~2 dowodu twierdzenia~\ref{thm:photon-emergence} wspomniano:
``$\oint\partial_\mu\theta\,dl^\mu = 2\pi n$, $n\in\mathbb{Z}$''.
\textbf{Problem}: nie podano, dlaczego $n$ jest całkowite,
dlaczego ładunek elektryczny jest kwantowany, ani jak
połączyć topologię wirów z~dyskretnym spektrum ładunku.

\subsubsection*{O.3.2 Kompaktowość fazy substratu}

\begin{lemat}[Kompaktowość fazy lokalnej]
\label{lem:compact_phase}
Faza lokalna $\theta_i \in [0, 2\pi)$ jest zmienną kompaktową
(periodicznie zidentyfikowaną: $\theta_i \equiv \theta_i + 2\pi$)
z~następujących przyczyn:
\begin{enumerate}
  \item Hamiltoniano~\eqref{eq:H-complex} zależy od $\theta_i$
        wyłącznie przez $\cos(\theta_j - \theta_i)$, który jest
        $2\pi$-periodyczny.
  \item Stan substratowy $\psi_i = \phi_i e^{i\theta_i}$ jest
        fizycznie tożsamy dla $\theta_i$ i~$\theta_i + 2\pi$
        (niezmienniczość obserwabli: $|\psi_i|$, $\mathrm{Re}(\psi_i^*\psi_j)$).
  \item Miara całkowania $\int_0^{2\pi} d\theta_i$ jest naturalną
        miarą Haara grupy $U(1)$.
\end{enumerate}
Zatem przestrzeń konfiguracyjna faz jest produktem $[0,2\pi)^N$,
topologicznie $T^N$ (N-wymiarowy torus).
\end{lemat}

\subsubsection*{O.3.3 Wiry fazowe jako topologiczne defekty}

\begin{definicja}[Wir fazowy]
\label{def:vortex}
Wir fazowy z~liczbą zwojową $n \in \mathbb{Z}$ jest konfiguracją
$\{\theta_i\}_{i\in V}$ taką, że dla dowolnego cyklu
$\gamma = (i_1, i_2, \ldots, i_k, i_1)$ otaczającego
punkt defektu $x_0$:
\begin{equation}
  \sum_{l=1}^k \Delta\theta_{i_l i_{l+1}} = 2\pi n,
  \qquad
  \Delta\theta_{ij} \equiv \mathrm{Im}\!\left[\ln\frac{\psi_j}{\psi_i}\right]
  \in (-\pi, \pi],
  \label{eq:winding_def}
\end{equation}
gdzie gałąź logarytmu dobrana tak, by różnica faz sąsiednich
węzłów leżała w~$(-\pi,\pi]$.
\end{definicja}

\begin{twierdzenie}[Kwantyzacja liczby zwojowej]
\label{thm:winding_quant}
Liczba zwojowa $n[\gamma]$ jest całkowita dla każdej pętli $\gamma$.
\end{twierdzenie}

\begin{proof}
Z~kompaktowości fazy ($\theta_i \in [0,2\pi)$, lemat~\ref{lem:compact_phase}):
skumulowana zmiana fazy wzdłuż zamkniętej pętli wynosi:
\begin{equation}
  \oint_\gamma d\theta
  = \theta_\mathrm{final} - \theta_\mathrm{initial}
  = 2\pi n, \quad n \in \mathbb{Z},
\end{equation}
bo $\theta_\mathrm{final} \equiv \theta_\mathrm{initial}$
(jednoznaczność stanu po pełnym obiegu), a~różnica musi być
wielokrotnością okresu $2\pi$.
\end{proof}

\subsubsection*{O.3.4 Kwantyzacja ładunku elektrycznego}

\begin{twierdzenie}[Kwantyzacja ładunku z topologii]
\label{thm:charge_quantization}
Ładunek elektryczny sprzężony z~polem fazowym $\theta_i$ jest
kwantowany: $q = n \cdot e_0$, $n \in \mathbb{Z}$, gdzie
$e_0$ jest ładunkiem elementarnym.
\end{twierdzenie}

\begin{proof}
Rozważamy pole materii substratowej sprzężone z~fazą:
pole skalarne $\Psi_\mathrm{matter}(x)$ o~ładunku $q$ transformuje się jak:
\begin{equation}
  \Psi_\mathrm{matter}(x) \to e^{iq\lambda(x)/\hbar_0}\,\Psi_\mathrm{matter}(x)
  \quad\text{przy}\quad
  \theta(x) \to \theta(x) + \lambda(x).
  \label{eq:matter_coupling}
\end{equation}
Jednoznaczność $\Psi_\mathrm{matter}$ wymaga, by po obejściu wokół
wiru z~$n[\gamma] = 1$:
\begin{equation}
  \exp\!\left(\frac{iq}{\hbar_0}\oint_\gamma d\lambda\right)
  = \exp\!\left(\frac{iq}{\hbar_0}\cdot 2\pi\right)
  \stackrel{!}{=} 1,
\end{equation}
skąd: $q/\hbar_0 \in \mathbb{Z}$, czyli $q = n\hbar_0 \equiv n\cdot e_0$
dla $e_0 \equiv \hbar_0$ (w~jednostkach Plancka).

\textbf{Identyfikacja $e_0$:} Ładunek elementarny w~jednostkach SI:
\begin{equation}
  \boxed{
    e_0
    = \hbar_0 \cdot
      \frac{2\pi}{\Phi_\mathrm{mag,min}},
  }
  \label{eq:e0_quant}
\end{equation}
gdzie $\Phi_\mathrm{mag,min} = h_0/e_0$ jest minimalnym kvantem
strumienia magnetycznego --- dokładna analogia warunku Diraca
$e_0 g_0 = 2\pi\hbar_0$ (z~monopolem Diraca zastąpionym worem
fazowym substratu).
\end{proof}

\begin{uwaga}[Kwantyzacja Diraca z substratu]
Warunek kwantyzacji~\eqref{eq:e0_quant} jest substratową
\emph{realizacją} argumentu Diraca: tam monopol magnetyczny
jest założeniem, tu jest \textbf{nieuniknionym defektem topologicznym}
sieci o~kompaktowej fazie. Nie ma potrzeby postulowania monopoli
jako osobnych cząstek --- każdy wir w~substracie fazowym
wymusza kwantyzację $q$.
\end{uwaga}

\begin{uwaga}[Status Luki O-2]
Luka O-2 jest \textbf{zamknięta}. Kwantyzacja ładunku wynika
z~kompaktowości $\theta_i \in [0,2\pi)$ (lemat~\ref{lem:compact_phase})
i~twierdzenia~\ref{thm:winding_quant}: $n[\gamma] \in \mathbb{Z}$.
Warunek jednojednoznaczności pola materii~\eqref{eq:matter_coupling}
wymusza $q = n e_0$. Dowód jest topologiczny --- nie zależy od
szczegółów hamiltonianu.
\end{uwaga}

% ---------------------------------------------------------------------------
\subsection*{O.4 Luka O-3: Stała struktury subtelnej z substratu}
\label{ssec:alpha_em_derivation}

\subsubsection*{O.4.1 Problem}

Propozycja~\ref{prop:alpha-em} w~sek09 podaje wzór:
\begin{equation}
  \alpha_\mathrm{em}
  = \frac{J^2 a_\mathrm{sub}^4 v^4}
         {4\pi\hbar c_0\varepsilon_0^{-1}}.
  \label{eq:alpha_sek09}
\end{equation}
\textbf{Problem}: brak formalnego wyprowadzenia, niejasne
jest, skąd pochodzi czynnik $J^2$ (nie $J$), a~relacja
do~eq.~\eqref{eq:mu0-substrate} nie jest podana.

\subsubsection*{O.4.2 Wyprowadzenie z relacji $\mu_0$--substrat}

Z~twierdzenia~\ref{thm:photon-emergence} (Krok~3, eq.~\eqref{eq:mu0-substrate}):
\begin{equation}
  \frac{1}{\mu_0} = \frac{2Jv^2 a_\mathrm{sub}^2 e^2}{\hbar_0^2},
  \qquad \text{skąd:} \quad
  e^2 = \frac{\hbar_0^2}{2\mu_0 J v^2 a_\mathrm{sub}^2}.
  \label{eq:e2_from_mu0}
\end{equation}

Stała struktury subtelnej w~jednostkach Gaussa ($\alpha_\mathrm{em} = e^2/(\hbar_0 c_0)$,
pamiętając o~$\varepsilon_0 = 1/(c_0^2\mu_0)$):
\begin{align}
  \alpha_\mathrm{em}
  &= \frac{e^2}{4\pi\varepsilon_0\,\hbar_0 c_0}
  = \frac{e^2 \mu_0 c_0}{4\pi\hbar_0}
  \label{eq:alpha_def_SI}\\
  &= \frac{\hbar_0^2}{2\mu_0 J v^2 a_\mathrm{sub}^2}
     \cdot \frac{\mu_0 c_0}{4\pi\hbar_0}
  = \frac{\hbar_0 c_0}{8\pi J v^2 a_\mathrm{sub}^2}.
  \label{eq:alpha_derived}
\end{align}

\begin{propozycja}[Stała $\alpha_\mathrm{em}$ z substratu — wersja poprawiona]
\label{prop:alpha_em_formal}
W~jednostkach Plancka ($\hbar_0 = c_0 = 1$, $a_\mathrm{sub}$ w~jednostkach $\ell_P$):
\begin{equation}
  \boxed{
    \alpha_\mathrm{em}^\mathrm{(sub)}
    = \frac{1}{8\pi\, J\, v^2\, a_\mathrm{sub}^2},
  }
  \label{eq:alpha_formal}
\end{equation}
gdzie $v^2 = \langle\phi^2\rangle$ w~fazie substratowej i~$a_\mathrm{sub}$
jest krokiem sieci w~jednostkach długości Plancka.
\end{propozycja}

\begin{uwaga}[Rozbieżność z eq.~\eqref{eq:alpha_sek09}]
\label{rem:alpha_discrepancy}
Wzór~\eqref{eq:alpha_formal} ma $J^1$ zamiast $J^2$ w~\eqref{eq:alpha_sek09}.
Rozbieżność pochodzi z~tego, że w~\eqref{eq:alpha_sek09} użyto
$e^2 \sim Jv^2 a_\mathrm{sub}^2$ (z~$1/\mu_0$), a~następnie
$\alpha_\mathrm{em} \sim e^2/(e^2/J) = J$ przez nieświadome
podwójne użycie relacji. Wzór~\eqref{eq:alpha_formal} jest
\textbf{poprawny} jako wynikający bezpośrednio z~\eqref{eq:mu0-substrate}.
\end{uwaga}

\subsubsection*{O.4.3 Relacja do hipotezy $a_\Gamma \cdot \Phi_0 = 1$}

Identyfikując $v^2 = \langle\phi^2\rangle = \Phi_0$ (tło próżni, sek10)
oraz $a_\mathrm{sub} = a_\Gamma$ z~hipotezą samospójności $a_\Gamma\Phi_0 = 1$:
\begin{equation}
  a_\mathrm{sub} = a_\Gamma = \frac{1}{\Phi_0}
  \quad\Longrightarrow\quad
  a_\mathrm{sub}^2 = \frac{1}{\Phi_0^2}.
  \label{eq:asub_from_phi0}
\end{equation}

Podstawiając do~\eqref{eq:alpha_formal}:
\begin{equation}
  \alpha_\mathrm{em}^\mathrm{(sub)}
  = \frac{1}{8\pi J \cdot \Phi_0 \cdot \Phi_0^{-2}}
  = \frac{\Phi_0}{8\pi J},
  \label{eq:alpha_UV}
\end{equation}
gdzie stała sprzężenia substratowego $J$ jest \emph{jedynym} wolnym
parametrem sektora elektromagnetycznego (przy ustalonej hipotezie
$a_\Gamma\Phi_0 = 1$ i~$v^2 = \Phi_0$).

\begin{propozycja}[$J_\mathrm{sub}$ z warunku $\alpha_\mathrm{em}^\mathrm{(sub)}$]
\label{prop:J_from_alpha}
Wymagając, by $\alpha_\mathrm{em}^\mathrm{(sub)} = \alpha_\mathrm{em}(\ell_P)$
(stała sprzężenia wyznaczona przez przepływ RG, eq.~\eqref{eq:alpha-Planck}
w~sek09: $\alpha_\mathrm{em}(\ell_P) \approx 1/94.09$):
\begin{equation}
  J_\mathrm{sub}
  = \frac{\Phi_0 \cdot 94.09}{8\pi}
  = \frac{24.66 \times 94.09}{8\pi}
  \approx 92.5.
  \label{eq:J_from_alpha}
\end{equation}
\textbf{Uwaga:} Wartość $J_\mathrm{sub} \approx 92.5$ dotyczy
sprzężenia sektora \emph{fazowego} substratu (sektor U(1)).
Sprzężenie sektora \emph{amplitudowego} (które daje $T_c$ w~MC)
jest osobnym parametrem: oba mogą być rzędu jedności
w~różnych skalach energii.
\end{propozycja}

\begin{uwaga}[Dwa parametry sprzężenia a ontologia TGP]
TGP ma \emph{jeden} substrat $\Gamma = (V, E)$, ale dwa
\emph{sektory} sprzężeń:
\begin{itemize}
  \item $J_\mathrm{amp}$: sprzężenie amplitudowe ($|\psi_i|$) ---
        daje pole $\Phi$, geometrię, grawitację.
  \item $J_\mathrm{phase}$: sprzężenie fazowe ($\theta_i$) ---
        daje pole $A_\mu$, elektromagnetyzm.
\end{itemize}
Oba sprzężenia wchodzą do tego samego hamiltonianu~\eqref{eq:H-complex}
przez człon $-J\,\mathrm{Re}(\psi_i^*\psi_j) = -J\phi_i\phi_j\cos(\theta_j-\theta_i)$:
jest tu \emph{iloczyn} amplitud i~fazy. W~LPA (dodatek N):
przepływ RG może rozdzielić efektywne sprzężenia amplitudowe
i~fazowe nawet jeśli mikroskopicznie są równe $J_\mathrm{amp} = J_\mathrm{phase} = J$.
Wymaganie $J_\mathrm{amp} = J_\mathrm{phase}$ jest \emph{warunkiem},
nie koniecznością --- stanowi dodatkowy program badań.
\end{uwaga}

\subsubsection*{O.4.4 Weryfikacja konsystencji wzoru}

\begin{propozycja}[Relacja między $\alpha_\mathrm{em}$ i $\kappa$]
\label{prop:alpha_kappa}
Stała unifikacji TGP $\kappa = 3/(4\Phi_0)$ (sek08) i~stała
elektromagnetyczna są powiązane przez:
\begin{equation}
  \frac{\alpha_\mathrm{em}^\mathrm{(sub)}}{\kappa}
  = \frac{\Phi_0/(8\pi J)}{3/(4\Phi_0)}
  = \frac{\Phi_0^2}{6\pi J}.
  \label{eq:alpha_kappa_ratio}
\end{equation}
Dla $\Phi_0 = 24.66$, $J = J_\mathrm{phase}$:
\begin{equation}
  \frac{\alpha_\mathrm{em}^\mathrm{(sub)}}{\kappa}
  = \frac{24.66^2}{6\pi J_\mathrm{phase}}
  \approx \frac{607.7}{18.85\,J_\mathrm{phase}}
  \approx \frac{32.2}{J_\mathrm{phase}}.
\end{equation}
Dla $\alpha_\mathrm{em}^\mathrm{(sub)}/\kappa = 1$ (hipotetyczna unifikacja
EM z~grawitacją TGP na skali Plancka): $J_\mathrm{phase} \approx 32.2$.
Wartość $J_\mathrm{phase} \approx 32$ nie jest sprzeczna z~porządkiem
jedności w~jednostkach Plancka (sprzężenie silne UV).
Ta relacja \textbf{otwiera} nowy program badań: unifikacja EM
i~grawitacji TGP jako dwóch przepływów tego samego sprzężenia $J$
w~ERG.
\end{propozycja}

% ---------------------------------------------------------------------------
\subsection*{O.5 Rozkład pola $A_\mu$ a stopnie swobody fotonu}
\label{ssec:dof_photon}

\subsubsection*{O.5.1 Liczba fizycznych stopni swobody}

\begin{propozycja}[Dwa stopnie swobody fotonu w TGP]
\label{prop:2_dof_photon}
Pole $A_\mu(x)$ ma 4 składowe. Po nałożeniu:
\begin{enumerate}
  \item Warunku transwersalności $k^\mu \tilde{A}_\mu^T = 0$: eliminuje 1 DOF,
  \item Swobody cechowania $A_\mu \to A_\mu + \partial_\mu\lambda$: eliminuje 1 DOF (podłużne $A_\mu^L$),
\end{enumerate}
pozostają \textbf{2 niezależne fizyczne stopnie swobody}: dwie
polaryzacje transwersalne fotonu.

W~TGP stopnie te emergują bezpośrednio z~gradientu fazy:
$A_\mu^T = \Pi_{\mu\nu}^T(\partial^\nu\theta)$,
gdzie $\Pi_{\mu\nu}^T = \delta_{\mu\nu} - \partial_\mu\partial_\nu/\Box$
jest projektorem transwersalnym (operatorowo).
\end{propozycja}

\subsubsection*{O.5.2 Bezmasowość w granicy termodynamicznej}

\begin{twierdzenie}[Ochrona bezmasowości fotonu w TGP]
\label{thm:photon_massless_protection}
Foton emergentny $A_\mu^T$ pozostaje bezmasowy w~granicy
termodynamicznej ($N = |V| \to \infty$). Ochrona jest \emph{topologiczna}:
\begin{enumerate}
  \item Symetria $U(1)$: $\theta_i \to \theta_i + \alpha$ (globalnie)
        jest \textbf{dokładna} dla każdego $N$ (hamiltoniano~\eqref{eq:H-complex}
        zależy wyłącznie od różnic $\theta_j - \theta_i$).
  \item Goldstone: łamanie globalne $U(1) \to \{1\}$ (przy $T < T_c$)
        daje \emph{jeden} bozon Goldstone'a --- tryb fazowy $\delta\theta$.
  \item Identyfikacja: tryb Goldstone'a = foton ($A_\mu \propto \partial_\mu\theta$,
        prędkość $v_\mathrm{ph} = c_0$, masa $= 0$).
\end{enumerate}
Masa fotonu jest chroniona twierdzeniem Goldstone'a --- nie przez
fine-tuning ani przez zewnętrzną symetrię cechowania.
\end{twierdzenie}

% ---------------------------------------------------------------------------
\subsection*{O.6 Mapa zamknięcia luk}

\begin{center}
\begin{tabular}{llll}
\toprule
Luka & Krok w sek09 & Narzędzie & Status \\
\midrule
O-1 & Krok~3, $\supset$ & Lem.~\ref{lem:F_transverse}, Tw.~\ref{thm:supset_exact} & \textbf{[ZAMKNIĘTA]} \\
O-2 & Krok~2, wiry & Lem.~\ref{lem:compact_phase}, Tw.~\ref{thm:charge_quantization} & \textbf{[ZAMKNIĘTA]} \\
O-3 & Prop.~\ref{prop:alpha-em} & Prop.~\ref{prop:alpha_em_formal}, eq.~\eqref{eq:alpha_formal} & \textbf{[ZAMKNIĘTA]} \\
\bottomrule
\end{tabular}
\end{center}

\medskip\noindent
\textbf{Otwarte po domknięciu luk:}
\begin{itemize}
  \item $J_\mathrm{amp} \neq J_\mathrm{phase}$: czy przepływ ERG rozdziela oba sprzężenia?
        (program badań w~ramach LPA', dodatek N).
  \item $\alpha_\mathrm{em}$ z~parametrów $(J, \Phi_0)$ numerycznie:
        weryfikacja eq.~\eqref{eq:alpha_UV} skryptem \texttt{tgp\_agamma\_phi0\_test.py}.
  \item Formalne twierdzenie Goldstone'a dla sieci skończonej → granica $N\to\infty$.
\end{itemize}

\bigskip
\begin{center}
\fbox{\parbox{0.85\textwidth}{
\textbf{Podsumowanie Dodatku O.}
Trzy luki w~dowodzie emergencji U(1) (sek09) zostały zamknięte
bez postulowania nowych bytów:
(O-1) Rozkład $A_\mu = A_\mu^T + A_\mu^L$ jest formalną tożsamością
Helmholtza --- człon podłużny jest czystą cechowaniem.
(O-2) Kwantyzacja ładunku wynika z~kompaktowości $\theta_i \in [0,2\pi)$
i~twierdzenia o~wirach --- jest obowiązkiem topologicznym, nie
postulatem.
(O-3) Stała $\alpha_\mathrm{em} = \hbar_0 c_0/(8\pi J v^2 a_\mathrm{sub}^2)$
wynika bezpośrednio z~relacji $\mu_0$--substrat~\eqref{eq:mu0-substrate}.
}}
\end{center}
